\section{The equality locus}\label{sec:ties}

The inequality \eqref{eq:dominate} is an equality on a set of full
dimension.  At every size $m\ge k+1$ and every weight vector there is a
design whose best $k$-subset has $\lmin(\SC)=1$, and these designs form
families of dimension $m-1$, the dimension of the weights themselves.  No
argument for $\Gtz$ can therefore carry a margin, and none of the
reductions of \S\ref{sec:real} may consume one.

\subsection{Extremal designs}\label{sec:extremal}

\begin{definition}\label{def:tie}
A design is \emph{extremal} if some $k$-subset dominates and no $k$-subset
dominates strictly: there is a $k$-subset with $\SC\psd I_k$, and no
$k$-subset has $\SC\pd I_k$.
\end{definition}

For a Hermitian matrix $A$, the relation $A\psd I_k$ says
$\lmin(A)\ge1$ and $A\pd I_k$ says $\lmin(A)>1$.

\begin{lemma}\label{lem:tiemax}
A design is extremal if and only if
$\max_{\lvert C\rvert=k}\lmin(\SC)=1$.  At an extremal design every
dominating $k$-subset has $\lmin(\SC)=1$, so that $\SC-I_k$ is
positive semidefinite and singular.
\end{lemma}

\begin{proof}
The maximum is at least $1$ if and only if some $C$ dominates, and exceeds
$1$ if and only if some $C$ dominates strictly.  A dominating $C$ has
$\lmin(\SC)\ge1$, and the maximum being $1$ forces equality.
\end{proof}

For $m\ge k$ put
\begin{equation}\label{eq:alpha}
  \alpha(m,k)\;:=\;\inf_{\dsn}\ \max_{\lvert C\rvert=k}\lmin(\SC),
\end{equation}
the infimum over real designs of size $m$ and rank $k$; this is the real
counterpart of the constant of Remark~\ref{rem:sharpc}, refined by the
size.  By construction $\alpha_k(\R)=\inf_{m\ge k}\alpha(m,k)$, so the
assertion $\alpha_k(\R)\ge1$ of \S\ref{sec:cvalue} is $\Gtz(m,k)$ for
every $m$.  Since a maximum is at least the infimum of maxima,
$\Gtz(m,k)$ holds if and only if $\alpha(m,k)\ge1$, and a design is
extremal exactly when it realizes the value $1$.

At corank zero there is nothing to construct.  If $m=k$ and $k\ge2$ then
the unique $k$-subset is $\{1,\dots,m\}$, and \eqref{eq:coparseval} gives
\[
  \SC-I_k\;=\;\sum_{c}(1-t_c)\,g_c^{\vphantom{\mathsf{H}}}g_c\adj\;\pd\;0 ,
\]
because the $g_c$ are then a basis of $\FF^k$ and $m\ge2$ forces $t_c<1$
for every $c$.  So no design at $m=k$, $k\ge2$, is extremal;
Proposition~\ref{prop:square} computes its value, $1/\max_ct_c$.  Extremal
designs first occur at corank one, where by
Theorem~\ref{thm:corankone}(iii)--(iv) they occur at every weight vector.

\subsection{Splitting, and existence at every size}\label{sec:split}

Evaluating \eqref{eq:dominate} at a vector gives the \emph{gap form}: for
every $k$-subset $C$ and every $x\in\FF^k$,
\begin{equation}\label{eq:gapform}
  x\adj\bigl(\SC-I_k\bigr)x
  \;=\;\sum_{c\in C}\lvert\langle g_c,x\rangle\rvert^{2}
       \;-\;\lVert x\rVert^{2}.
\end{equation}
Neither side involves the weights.

\begin{lemma}\label{lem:parallel}
Let $C$ be a $k$-subset containing two distinct indices $c,c'$ with
$g_{c'}=a\,g_c$ for some scalar $a$.  Then $\SC$ is singular and $C$
does not dominate.
\end{lemma}

\begin{proof}
The family $(g_c)_{c\in C}$ has $k$ members and spans a subspace of
dimension at most $k-1$, so some unit vector $x$ is orthogonal to all of
them.  Then $\SC x=0$, and by \eqref{eq:gapform},
$x\adj(\SC-I_k)x=-1<0$.
\end{proof}

Now let $\dsn$ be a design at $(m,k)$, let $e$ be an index and let
$0<a<t_e$.  Replace the pair $(g_e,t_e)$ by the two pairs $(g_e,a)$ and
$(g_e,t_e-a)$, leaving everything else alone.  The weights remain positive
and still sum to one, and
\[
  a\,g_e^{\vphantom{\mathsf{H}}}g_e\adj+(t_e-a)\,g_e^{\vphantom{\mathsf{H}}}g_e\adj
  \;=\;t_e\,g_e^{\vphantom{\mathsf{H}}}g_e\adj,
\]
so \eqref{eq:parseval} is unchanged and the result is a design at
$(m+1,k)$.  The vector is not rescaled: \eqref{eq:parseval} is affected
only through the total weight carried by a direction.  This is the
\emph{split} at $e$ with ratio $a$.

Iterating splits from a design at $(k+1,k)$ produces, at size $m$, a
design whose $m$ vectors take at most $k+1$ distinct values and whose
weights are arbitrary subject to prescribed class totals.  We construct
these directly.

\begin{theorem}[Lean]\label{thm:tieseverywhere}
Let $k\ge1$ and $m\ge k+1$, let $t$ be a weight vector on $m$ atoms, and
let $\varphi:\{1,\dots,m\}\to\{1,\dots,k+1\}$ be a surjection.  Put
\[
  T_j\;:=\;\sum_{\varphi(c)=j}t_c \qquad (1\le j\le k+1),
\]
let $(h_j,T_j)_{j\le k+1}$ be a design at $(k+1,k)$ with weights $T$ in
which every $k$-subset dominates and none dominates strictly, and set
$g_c:=h_{\varphi(c)}$.  Then $(g_c,t_c)_{c\le m}$ is an extremal design at
$(m,k)$ with weights exactly $t$.  In particular extremal designs exist at
$(m,k)$ for every $m\ge k+1$, over either field, and at every weight
vector.
\end{theorem}

\begin{proof}
The $T_j$ are positive and sum to one, so Theorem~\ref{thm:corankone}(iii)--(iv)
supplies the design $(h_j,T_j)$ demanded.  Grouping the sum
\eqref{eq:parseval} by classes,
\[
  \sum_{c=1}^{m}t_c\,g_c^{\vphantom{\mathsf{H}}}g_c\adj
  \;=\;\sum_{j=1}^{k+1}\Bigl(\sum_{\varphi(c)=j}t_c\Bigr)
        h_j^{\vphantom{\mathsf{H}}}h_j\adj
  \;=\;\sum_{j=1}^{k+1}T_j\,h_j^{\vphantom{\mathsf{H}}}h_j\adj\;=\;I_k,
\]
so $(g,t)$ is a design at $(m,k)$, with weights $t$ by construction.

Let $C$ be a $k$-subset.  If $\varphi$ is not injective on $C$ then $C$
carries two equal vectors and does not dominate, by
Lemma~\ref{lem:parallel}.  If $\varphi$ is injective on $C$ then
$\varphi(C)$ is a $k$-subset of $\{1,\dots,k+1\}$ and
\[
  \SC\;=\;\sum_{c\in C}h_{\varphi(c)}^{\vphantom{\mathsf{H}}}h_{\varphi(c)}\adj
      \;=\;S_{\varphi(C)},
\]
which dominates and does not dominate strictly.  Since $\varphi$ is
surjective there is at least one $C$ of the second kind: choose $k$ of the
$k+1$ classes and one index from each.  Hence some $k$-subset dominates
and none dominates strictly.
\end{proof}

At $m=k+1$ the surjection is a bijection and the theorem returns
Theorem~\ref{thm:corankone}(iii)--(iv).  At $m\ge k+2$ some class contains two indices,
so the design carries a pair of equal vectors; \S\ref{sec:diamond} shows
that this feature is not forced.

\begin{proposition}\label{prop:section}
For fixed $\varphi$ the assignment $t\mapsto \dsn(t,\varphi)$ of
Theorem~\ref{thm:tieseverywhere} is a section of the weight map over the
whole open simplex: every weight vector carries an extremal design.  At
$m=k+1$ it is onto, in the following sense: two extremal designs of size
$k+1$ with the same weights differ by
\[
  g_c\;\longmapsto\;\varepsilon_c\,Ug_c,
  \qquad U\in O(k)\ \text{resp.}\ U(k),\qquad
  \lvert\varepsilon_c\rvert=1 ,
\]
a transformation under which every $\SC$ becomes $U\SC U\adj$, so that no
$\lmin(\SC)$ changes and neither domination nor extremality is
affected.  Consequently the extremal
locus at corank one has dimension $m-1=k$ modulo that group, and
at every $m\ge k+1$ the extremal locus contains, for each partition of
$\{1,\dots,m\}$ into $k+1$ nonempty blocks, a family of dimension $m-1$.
\end{proposition}

\begin{proof}
Only the corank-one assertion needs an argument.  Let $m=k+1$.  By
Theorem~\ref{thm:corankone} a design of that size with weights $t$ is
extremal if and only if its chart projection is $P=I-zz\adj$ with
$\lvert z_c\rvert^{2}=(1-t_c)/k$ for every $c$.  One direction is part
(iii) of that theorem.  For the other, extremality forbids strict
domination, so by part (i) no $p_{c_0}$ exceeds the $t$-mean
$\sum_ct_cp_c$; the weights being positive, a family none of whose members
exceeds its own weighted mean is constant, and by \eqref{eq:pnorm} the
constant is $1/k$, which is the stated condition on $z$.  These moduli
determine $z$ up to the individual phases $\varepsilon_c$, and $P$
determines $V$ up to $V\mapsto VU$, since two isometries with the same
range differ by the unitary $U=V\adj V'$.  Replacing $z$ by
$\diag(\varepsilon)z$ replaces $P$ by
$\diag(\varepsilon)P\diag(\varepsilon)\adj$, the identity term surviving
that conjugation because $\lvert\varepsilon_c\rvert=1$; so $V$ becomes
$\diag(\varepsilon)V$ and $g_c=v_c/\sqrt{t_c}$ becomes $\varepsilon_cg_c$,
under which $g_c^{\vphantom{\mathsf{H}}}g_c\adj$ is unchanged.  The fibre of
the weight map over each $t$ is therefore a single orbit, the
quotient is mapped bijectively onto the open simplex, and the dimension is
$m-1=k$.
\end{proof}

The corank-one classification does not extend: the equations
\eqref{eq:tielaw} that cut it out have no solutions above corank one.

\begin{proposition}[Lean]\label{prop:tielawcorank}
Let $k\ge2$.  If a design of size $m$ satisfies $k\,t_c\ell_c=k-1+t_c$ at
every atom, then $m=k+1$.
\end{proposition}

\begin{proof}
Sum \eqref{eq:tielaw} over $c$.  The left side is
$k\sum_ct_c\ell_c=k^{2}$ by Lemma~\ref{lem:trace}; the right side is
$m(k-1)+\sum_ct_c=m(k-1)+1$.  Hence $(k-1)(k+1)=m(k-1)$, and $k\ge2$
permits the cancellation.
\end{proof}

So above corank one the extremal designs are cut out by no such closed
system.  The designs of Theorem~\ref{thm:tieseverywhere} satisfy
\eqref{eq:tielaw} at an atom $c$ precisely when $c$ is alone in its class:
substituting $\ell_c=(k-1+T_{\varphi(c)})/(kT_{\varphi(c)})$, which is
Theorem~\ref{thm:corankone}(iii) read at the class totals, into
\eqref{eq:tielaw} cancels the products and leaves
$t_c(k-1)=T_{\varphi(c)}(k-1)$, that is $t_c=T_{\varphi(c)}$.  The law
therefore fails at exactly the atoms of the classes that carry more than
one, and at $m=k+2$ it fails at two atoms out of $k+2$.  The following
instances are used again in \S\ref{sec:interior} and in
Appendix~\ref{app:data}.

\begin{example}[Lean]\label{ex:split}
Take for $(h_j)$ the regular tetrahedron of Example~\ref{ex:tetra}, so
that $k=3$, $T_j=\tfrac14$, $\ell(h_j)=3$ and
$\langle h_i,h_j\rangle=-1$ for $i\ne j$.

At $m=6$, with $\varphi=(1,2,3,3,4,4)$ and weights
$\bigl(\tfrac14,\tfrac14,a,\tfrac14-a,b,\tfrac14-b\bigr)$ for any
$0<a,b<\tfrac14$, Theorem~\ref{thm:tieseverywhere} produces a
two-parameter family of extremal designs at $(6,3)$.  Every leverage is
$3$, while \eqref{eq:tielaw} at $k=3$ reads $9t_c=2+t_c$, that is
$t_c=\tfrac14$: the law fails at each of the four split atoms, in
accordance with Proposition~\ref{prop:tielawcorank}.

At $m=7$, with $\varphi=(1,2,3,4,1,2,3)$ and weights
$\bigl(\tfrac18,\tfrac18,\tfrac18,\tfrac14,\tfrac18,\tfrac18,
\tfrac18\bigr)$, one obtains an extremal design at $(7,3)$ in which again
every leverage is $3$.  Of its $35$ triples, the $20$ carrying three
distinct classes dominate: for such a triple $T$, missing the class $j$,
the computation of Example~\ref{ex:tetra} gives
\[
  S_T\;=\;4I_3-h_j^{\vphantom{\mathsf{T}}}h_j\tp,
  \qquad
  x\tp(S_T-I_3)x\;=\;3\lVert x\rVert^{2}-\langle h_j,x\rangle^{2},
\]
so that the gap has spectrum $\{0,3,3\}$ and $\lmin(S_T)=1$.  The
remaining $15$ triples repeat a class and do not dominate, by
Lemma~\ref{lem:parallel}.
\end{example}

\subsection{No margin}\label{sec:nomargin}

\begin{corollary}[Lean]\label{cor:nomargin}
Let $k\ge1$, $m\ge k+1$ and $\varepsilon>0$.  For every weight vector $t$
there is a design at $(m,k)$ with weights $t$ admitting no $k$-subset with
$\SC\psd(1+\varepsilon)I_k$.  Consequently $\alpha(m,k)\le1$, and if
$\Gtz(m,k)$ holds then $\alpha(m,k)=1$ and the infimum in
\eqref{eq:alpha} is attained.
\end{corollary}

\begin{proof}
Let $\dsn$ be the extremal design of Theorem~\ref{thm:tieseverywhere} with
weights $t$.  If some $k$-subset had $\SC\psd(1+\varepsilon)I_k$ then
$\SC-I_k\psd\varepsilon I_k\pd0$, so that $C$ would dominate strictly.
The remaining assertions are Lemma~\ref{lem:tiemax}.
\end{proof}

There is no uniform margin.  For every $\varepsilon>0$ the assertion
``every design admits a $k$-subset with $\SC\psd(1+\varepsilon)I_k$'' is
false, at every cell with $m\ge k+1$.  An argument whose output is a
strict inequality --- a positivity certificate carrying a strictly
positive multiplier, or a bound $\lmin(\SC)\ge1+f$ with $f$ a positive
function of the design data --- therefore fails at the extremal designs
themselves, whatever its degree or complexity.

There is no slack to transport.  Suppose $\Gtz(m,k)$ is derived from a
smaller cell $(m',k')$ with $m'\ge k'+1$.  The slack
$\max_C\lmin(\SC)-1$ available at the smaller cell is zero at the designs
of Theorem~\ref{thm:tieseverywhere}, so no step may consume a positive
amount of it: the reduction must be an identity and not an estimate.  The
compression lift of \S\ref{sec:boundary} is such an identity at vanishing
weight, and Theorem~\ref{thm:sharp} exhibits its failure at every positive
weight, on the extremal locus itself.

The locus has full dimension.  By
Proposition~\ref{prop:section} the extremal locus meets every point of the
open weight simplex, at every size $m\ge k+1$, and has dimension $m-1$,
which is the dimension of the weights.  No restriction on the weights
avoids it, and no compactness argument on a subregion of the design space
yields a positive infimum for the slack unless the closure of that
subregion misses the locus.  The analysis at the extremal points must
therefore be first-order, and \S\ref{sec:compactness} and
\S\ref{sec:interior} supply it.  Appendix~\ref{app:graveyard} collects
the strategies that fail for these reasons, together with their refuting
witnesses.  Figure~\ref{fig:section} draws the locus at $(6,3)$.

\begin{figure}[ht]
  \centering
  % The equality locus of section sec:ties as a section of the weight map, at
% (m,k) = (6,3).  Orthographic projection of R^3 with azimuth 20 and elevation
% 15 degrees, scaled by 1.35cm:
%   X = 1.35(-0.3420 x_1 + 0.9397 x_2)
%   Y = 1.35(-0.2432 x_1 - 0.0885 x_2 + 0.9659 x_3)
% so one unit of height x_3 lifts a point by 1.35*0.9659 = 1.3040cm.
%
% FLOOR.  The base drawn is the two-dimensional slice t_4=t_5=t_6=1/6 of the
% open weight simplex at m=6; on it t_1+t_2+t_3=1/2, and the barycentric
% coordinate u=2(t_1,t_2,t_3) carries the slice to the equilateral triangle of
% circumradius 2.8 in the plane x_3=0,
%   V1 = ( 2.8, 0.0000, 0),  V2 = (-1.4, 2.4249, 0),  V3 = (-1.4,-2.4249, 0),
% the uniform weight vector t_c=1/6 landing on the centroid u=(1/3,1/3,1/3).
% Every coordinate below is u_1 V1 + u_2 V2 + u_3 V3 projected and then raised
% by 1.3040 z, with u and z in the trailing comment.
%
% SHEETS.  The fibre of the weight map is drawn as one vertical dimension and
% the three heights carry no numerical meaning.  Sheet j is the graph of
%   z = z0 + a q(u),   q(u) = u_1u_2 + u_2u_3 + u_3u_1,
%   (z0,a) = (1.60, 0.30), (3.05,-0.28), (4.50, 0.32),
% and q is 0 at the corners, 1/4 at the edge midpoints, 1/3 at the centroid.
% Each boundary edge is the quadratic Bezier through the corner, edge-midpoint
% and corner heights, written as the equal cubic: its two control points sit at
% the thirds of the chord, raised by 4d/3, where d = 1.3040 a/4 is the height of
% the edge midpoint above the chord midpoint,
%   d    = 0.0978, -0.0913, 0.1043,
%   4d/3 = 0.1304, -0.1217, 0.1391.
%
% A dashed boundary marks exclusion.  The simplex is open, and so is each sheet
% over it.  The subregion is the parallelogram of barycentric half-widths 0.05
% about u = (.36,.16,.48), drawn on the floor and on each of the three sheets.
\begin{tikzpicture}[
  x=1cm, y=1cm,
  vec/.style={-{Latex[length=2mm,width=1.5mm]}, line width=0.75pt},
  open/.style={draw=black!70, line width=0.45pt, dashed,
               dash pattern=on 2.2pt off 1.5pt},
  frame/.style={draw=black!35, line width=0.3pt, densely dotted},
  every node/.style={font=\small}]

% ---- the floor: the slice of the open weight simplex --------------------
\filldraw[fill=black!12, open]
  (-1.2928,-0.9193) --                              % u = (1,0,0)
  ( 3.7226, 0.1699) --                              % u = (0,1,0)
  (-2.4297, 0.7494) -- cycle;                       % u = (0,0,1)
\path[fill=black!35]
  (-0.8421, 0.1938) --                              % u = (.26,.21,.53)
  (-0.7284, 0.0270) --                              % u = (.36,.21,.43)
  (-1.2300,-0.0819) --                              % u = (.46,.11,.43)
  (-1.3437, 0.0849) -- cycle;                       % u = (.36,.11,.53)

% ---- sheet 1: the classes 1 | 2 | 35 | 46 -------------------------------
\filldraw[fill=black!12, open]
  (-1.2928,1.1671)                                  % u = (1,0,0), z = 1.60
    .. controls ( 0.3790,1.6606) and ( 2.0508,2.0236) ..
  ( 3.7226,2.2563)                                  % u = (0,1,0), z = 1.60
    .. controls ( 1.6718,2.5799) and (-0.3789,2.7730) ..
  (-2.4297,2.8358)                                  % u = (0,0,1), z = 1.60
    .. controls (-2.0507,2.4100) and (-1.6718,1.8537) .. cycle;
\path[fill=black!35]
  (-0.8421,2.3990) --                               % z = 1.6911
  (-0.7284,2.2389) --                               % z = 1.6962
  (-1.2300,2.1202) --                               % z = 1.6887
  (-1.3437,2.2842) -- cycle;                        % z = 1.6866

% ---- sheet 2: the classes 12 | 34 | 5 | 6 -------------------------------
\filldraw[fill=black!12, open]
  (-1.2928,3.0579)                                  % u = (1,0,0), z = 3.05
    .. controls ( 0.3790,3.2993) and ( 2.0508,3.6623) ..
  ( 3.7226,4.1471)                                  % u = (0,1,0), z = 3.05
    .. controls ( 1.6718,4.2186) and (-0.3789,4.4117) ..
  (-2.4297,4.7266)                                  % u = (0,0,1), z = 3.05
    .. controls (-2.0507,4.0487) and (-1.6718,3.4924) .. cycle;
\path[fill=black!35]
  (-0.8421,4.0601) --                               % z = 2.9650
  (-0.7284,3.8871) --                               % z = 2.9602
  (-1.2300,3.7873) --                               % z = 2.9672
  (-1.3437,3.9567) -- cycle;                        % z = 2.9692

% ---- sheet 3: the classes 156 | 2 | 3 | 4 -------------------------------
\filldraw[fill=black!12, open]
  (-1.2928,4.9487)                                  % u = (1,0,0), z = 4.50
    .. controls ( 0.3790,5.4509) and ( 2.0508,5.8139) ..
  ( 3.7226,6.0379)                                  % u = (0,1,0), z = 4.50
    .. controls ( 1.6718,6.3702) and (-0.3789,6.5633) ..
  (-2.4297,6.6174)                                  % u = (0,0,1), z = 4.50
    .. controls (-2.0507,6.2003) and (-1.6718,5.6440) .. cycle;
\path[fill=black!35]
  (-0.8421,6.1885) --                               % z = 4.5972
  (-0.7284,6.0288) --                               % z = 4.6026
  (-1.2300,5.9095) --                               % z = 4.5946
  (-1.3437,6.0734) -- cycle;                        % z = 4.5924

% ---- the fibre direction: three corner walls and two whole fibres -------
\draw[frame] (-1.2928,-0.9193) -- (-1.2928,5.05);
\draw[frame] ( 3.7226, 0.1699) -- ( 3.7226,6.14);
\draw[frame] (-2.4297, 0.7494) -- (-2.4297,6.72);
\draw[frame] ( 0.0000, 0.0000) -- ( 0.0000,6.70);   % u = (1/3,1/3,1/3)
\draw[frame] ( 1.9130, 0.1217) -- ( 1.9130,6.50);   % u = (.14,.68,.18)

% ---- the design each fibre meets on each sheet --------------------------
\fill ( 0.0000,0.0000) circle (1.6pt);
\fill ( 0.0000,2.2168) circle (1.6pt);              % z = 1.7000
\fill ( 0.0000,3.8555) circle (1.6pt);              % z = 2.9567
\fill ( 0.0000,6.0071) circle (1.6pt);              % z = 4.6067
\fill ( 1.9130,0.1217) circle (1.6pt);
\fill ( 1.9130,2.3031) circle (1.6pt);              % z = 1.6728
\fill ( 1.9130,4.0102) circle (1.6pt);              % z = 2.9820
\fill ( 1.9130,6.0910) circle (1.6pt);              % z = 4.5777

% ---- the section over the subregion -------------------------------------
\draw[vec] (-1.0361,0.0560) -- (-1.0361,2.1661);    % u = (.36,.16,.48)
\draw[line width=0.3pt] (-1.0361,1.02) -- (0.25,1.02);
\node[anchor=west, font=\scriptsize, align=left, fill=white, inner sep=1.5pt]
  at (0.29,1.02) {$t\mapsto\dsn(t,\varphi)$\\ at every $t\in\mathcal{U}$};
\node[anchor=north, font=\footnotesize] at (-1.05,-0.09) {$\mathcal{U}$};

% ---- the base, labelled -------------------------------------------------
\node[anchor=north east, font=\footnotesize] at (-1.36,-0.98) {$t_1=\tfrac12$};
\node[anchor=west,       font=\footnotesize] at ( 3.83,0.1699) {$t_2=\tfrac12$};
\node[anchor=east,       font=\footnotesize] at (-2.53,0.7494) {$t_3=\tfrac12$};
\node[anchor=west,  font=\scriptsize] at (0.12,-0.05) {$t_c=\tfrac16$};
\node[anchor=north, font=\scriptsize] at (0.90,-1.34)
  {$t_4=t_5=t_6=\tfrac16$, \quad $t_1+t_2+t_3=\tfrac12$};

% ---- the fibre axis and the classes -------------------------------------
\draw[vec, black!55] (-3.75,0.10) -- (-3.75,6.55);
\node[rotate=90, font=\scriptsize] at (-3.98,3.30) {fibre of the weight map};
\node[anchor=west, font=\scriptsize]   at (3.85,6.42) {classes of $\varphi$:};
\node[anchor=west, font=\footnotesize] at (3.85,6.0379) {$156\,|\,2\,|\,3\,|\,4$};
\node[anchor=west, font=\footnotesize] at (3.85,4.1471) {$12\,|\,34\,|\,5\,|\,6$};
\node[anchor=west, font=\footnotesize] at (3.85,2.2563) {$1\,|\,2\,|\,35\,|\,46$};

\end{tikzpicture}

  \caption{The equality locus as a section, at $(6,3)$.  Floor: the slice
  $t_4=t_5=t_6=\tfrac16$ of the open weight simplex, its excluded boundary
  dashed and the uniform point marked.  Vertical: the fibre of the weight
  map, drawn as a single dimension.  Each sheet is a section
  $t\mapsto\dsn(t,\varphi)$ of Proposition~\ref{prop:section} restricted to
  that slice; the section is defined at every weight vector and its image has
  dimension $m-1=5$, the dimension of the weights.  One sheet occurs for each
  partition of $\{1,\dots,6\}$ into four classes, $65$ in all, of which three
  are drawn.  However small, the shaded $\mathcal{U}$ meets every one of them,
  so no restriction of the weights isolates a region on which the slack has a
  positive infimum; four of the nine strategies of
  Appendix~\ref{app:graveyard} fail on that account.}
  \label{fig:section}
\end{figure}

\subsection{The diamond}\label{sec:diamond}

Every design of Theorem~\ref{thm:tieseverywhere} of size $m\ge k+2$
carries a pair of equal vectors and takes at most $k+1$ distinct
values; neither feature is forced.  We exhibit an extremal design at
$(5,3)$ whose five vectors are pairwise non-parallel.  It is described
by a graph, and its data are rational even though its vectors are not.

Let $\gph$ be the complete graph on the vertices $1,2,3,4$ with the edge
$34$ deleted, grounded at the vertex $4$, and list its five edges as
\[
  e_0=12,\qquad e_1=13,\qquad e_2=14,\qquad e_3=23,\qquad e_4=24 .
\]
A potential is a vector $y=(y_1,y_2,y_3)\in\R^3$, the ground carrying
$y_4=0$, and the \emph{drop} across an edge is the difference of the
potentials at its endpoints:
\begin{equation}\label{eq:drops}
  d(y)\;=\;\bigl(y_1-y_2,\ y_1-y_3,\ y_1,\ y_2-y_3,\ y_2\bigr)
  \;=\;\bigl(\langle b_e,y\rangle\bigr)_{0\le e\le4},
\end{equation}
the reduced incidence vectors being
\[
  b_0=(1,-1,0),\quad b_1=(1,0,-1),\quad b_2=(1,0,0),\quad
  b_3=(0,1,-1),\quad b_4=(0,1,0).
\]
Give the edges the conductances and weights
\[
  c\;=\;(2,3,3,3,3), \qquad t_e\;=\;\tfrac15\quad(0\le e\le4),
\]
and form the weighted Laplacian
\begin{equation}\label{eq:laplacian}
  L\;:=\;\sum_{e=0}^{4}c_e\,b_e^{\vphantom{\mathsf{T}}}b_e\tp
  \;=\;\begin{pmatrix} 8&-2&-3\\ -2&8&-3\\ -3&-3&6\end{pmatrix},
  \qquad
  y\tp Ly\;=\;\sum_{e}c_e\,d_e(y)^{2}.
\end{equation}
Its leading principal minors are $8$, $60$ and $\det L=180$, so $L\pd0$.
Choose any $R$ with $R\tp LR=I_3$ --- for instance $R=L^{-1/2}$ --- and
put
\begin{equation}\label{eq:diamondatoms}
  g_e\;:=\;\sqrt{5c_e}\;R\tp b_e \qquad (0\le e\le 4).
\end{equation}
Then
\[
  \sum_{e}t_e\,g_e^{\vphantom{\mathsf{T}}}g_e\tp
  \;=\;\sum_e \tfrac15\cdot 5c_e\,R\tp b_e^{\vphantom{\mathsf{T}}}b_e\tp R
  \;=\;R\tp LR\;=\;I_3,
\]
so $(g_e,t_e)$ is a design at $(5,3)$; call it $\dsn_{\Diamond}$.  It is
determined by the graph data up to the ambiguity $R\mapsto RU$ with $U$
orthogonal, which replaces every $g_e$ by $U\tp g_e$ and every $\SC$ by
$U\tp\SC U$, changing no $\lmin(\SC)$.

Since $RR\tp=L^{-1}$,
\begin{equation}\label{eq:diamondlev}
  \ell_e\;=\;5c_e\,b_e\tp L^{-1}b_e ,
  \qquad
  L^{-1}\;=\;\frac1{180}
  \begin{pmatrix}39&21&30\\ 21&39&30\\ 30&30&60\end{pmatrix},
\end{equation}
whence $b_0\tp L^{-1}b_0=\tfrac{36}{180}=\tfrac15$ and
$b_e\tp L^{-1}b_e=\tfrac{39}{180}=\tfrac{13}{60}$ for $e\ge1$, so that
\[
  \ell_0=2,\qquad \ell_1=\ell_2=\ell_3=\ell_4=\tfrac{13}4 .
\]
As a check, $\sum_et_e\ell_e=\tfrac15\bigl(2+4\cdot\tfrac{13}4\bigr)=3=k$,
as Lemma~\ref{lem:trace} requires.  The substitution $y=Rx$, a bijection
of $\R^3$, gives for every $3$-subset $C$
\[
  x\tp\SC x=\sum_{e\in C}5c_e\,d_e(y)^{2},
  \qquad
  x\tp x=x\tp R\tp LRx=\sum_{e}c_e\,d_e(y)^{2},
\]
so that domination becomes a statement about potentials alone, free of
the whitening $R$:
\begin{equation}\label{eq:diamondgap}
  x\tp\bigl(\SC-I_3\bigr)x\;=\;Q_C(y),
  \qquad
  Q_C(y):=\sum_{e=0}^{4}\bigl(5\cdot[e\in C]-1\bigr)c_e\,d_e(y)^{2},
\end{equation}
where $[\,\cdot\,]$ is $1$ when the enclosed statement holds and $0$
otherwise; and $C$ dominates if and only if $Q_C\ge0$ on $\R^3$, and
dominates strictly if and only if $Q_C>0$ off the origin.

\begin{example}[Lean]\label{ex:diamond}
The design $\dsn_{\Diamond}$ is extremal at $(5,3)$; its five vectors are
pairwise non-parallel; and it is not one of the designs of
Theorem~\ref{thm:tieseverywhere}.
\end{example}

\begin{proof}
For the star at the vertex $1$, that is $C=\{e_0,e_1,e_2\}$, expanding
\eqref{eq:diamondgap} by means of \eqref{eq:drops} gives
\begin{align*}
  Q_C(y)&=8(y_1-y_2)^2+12(y_1-y_3)^2+12y_1^2-3(y_2-y_3)^2-3y_2^2\\
        &=32y_1^2+2y_2^2+9y_3^2-16y_1y_2-24y_1y_3+6y_2y_3 .
\end{align*}
Doubling and completing the square in $y_1$ leaves $9y_3^{2}$:
\begin{equation}\label{eq:diamondstar}
  2\,Q_C(y)\;=\;\bigl(8y_1-2y_2-3y_3\bigr)^{2}+9\,y_3^{2} .
\end{equation}
So $C$ dominates.  The right side vanishes at the nonzero $y=(1,4,0)$, so
$Q_C$ is singular and $C$ does not dominate strictly.

There are ten $3$-subsets: eight spanning trees of $\gph$, and its two
triangles $\{e_0,e_1,e_3\}$ on the vertices $1,2,3$ and
$\{e_0,e_2,e_4\}$ on $1,2,4$.  For each, the potential listed below gives
the stated value of $Q_C$:
\[
\begin{array}{llr@{\qquad}llr}
\{e_0,e_1,e_2\} & (1,4,0)  & 0 &
\{e_0,e_1,e_3\} & (1,1,1)  & -6\\
\{e_0,e_1,e_4\} & (4,1,5)  & 0 &
\{e_0,e_2,e_3\} & (1,4,5)  & 0\\
\{e_0,e_2,e_4\} & (0,0,1)  & -6 &
\{e_0,e_3,e_4\} & (4,1,0)  & 0\\
\{e_1,e_2,e_3\} & (2,8,5)  & 0 &
\{e_1,e_2,e_4\} & (3,-3,5) & 0\\
\{e_1,e_3,e_4\} & (8,2,5)  & 0 &
\{e_2,e_3,e_4\} & (-3,3,5) & 0
\end{array}
\]
Each entry is one evaluation of \eqref{eq:diamondgap}.  At
$\{e_1,e_2,e_4\}$ and $y=(3,-3,5)$, for instance, the drops
\eqref{eq:drops} are $(6,-2,3,-8,-3)$ and
\[
  Q_C(y)=-2\cdot36+12\cdot4+12\cdot9-3\cdot64+12\cdot9=0 .
\]
For each triangle the potential is constant on the triangle's vertex set,
so its three drops vanish and $Q_C$ reduces to minus the Dirichlet
contribution of the two edges leaving that set.  Each of those two carries
conductance $3$ and a drop of modulus $1$, so $Q_C=-6$.  Hence no
$3$-subset has $Q_C>0$ off the origin, that is none dominates strictly,
while by the first paragraph one dominates: $\dsn_{\Diamond}$ is extremal.

Proportionality of $g_e$ and $g_f$ is proportionality of $b_e$ and $b_f$,
since $R\tp$ is invertible and the factors $\sqrt{5c_e}$ do not
vanish.  Proportional nonzero vectors have equal support, and the five
supports $\{1,2\},\{1,3\},\{1\},\{2,3\},\{2\}$ are distinct.  So no two
vectors are parallel; in particular $\dsn_{\Diamond}$ carries five distinct
directions, whereas a design of Theorem~\ref{thm:tieseverywhere} carries
at most $k+1=4$.
\end{proof}

\begin{corollary}[Lean]\label{cor:noparallel}
The assertion ``every extremal design contains two parallel vectors'' is
true for every design of Theorem~\ref{thm:tieseverywhere} of size
$m\ge k+2$, and false at $(5,3)$.
\end{corollary}

\begin{proof}
If $m\ge k+2$ then $\varphi$ is not injective, so two indices carry equal
vectors.  Example~\ref{ex:diamond} is the counterexample.
\end{proof}

An argument branching on the presence of a parallel pair therefore
requires a separate treatment of $\dsn_{\Diamond}$, and of every extremal
design not obtained by splitting a smaller one.  Whether the tetrahedron
and the diamond generate the whole rank-three extremal locus is
Problem~\ref{prob:structure}.  Figure~\ref{fig:diamond} shows the graph
and its incidence vectors.

\begin{figure}[ht]
  \centering
  % Left: the graph Gamma = K_4 minus the edge 34, grounded at the vertex 4;
% the line weight of an edge is proportional to its conductance.
% Right: the five reduced incidence vectors, drawn in R^3 by the orthographic
% projection with azimuth 20 and elevation -28 degrees, scaled by 1.5cm:
%   X = 1.5(-0.3420 y_1 + 0.9397 y_2)
%   Y = 1.5( 0.4411 y_1 + 0.1606 y_2 + 0.8830 y_3)
\begin{tikzpicture}[
  x=1cm, y=1cm,
  ver/.style={circle, draw, line width=0.5pt, fill=white, inner sep=0pt,
              minimum size=4.6mm, font=\small},
  spine/.style={line width=0.6pt},
  rim/.style={line width=0.9pt},
  vec/.style={-{Latex[length=1.9mm,width=1.4mm]}, line width=0.7pt},
  every node/.style={font=\small}]

% ================= the graph =============================================
\begin{scope}
  \coordinate (v1) at (-1.05, 0);
  \coordinate (v2) at ( 1.05, 0);
  \coordinate (v3) at ( 0, 1.15);
  \coordinate (v4) at ( 0,-1.15);

  \draw[rim]   (v1)--(v3) node[midway, above left=-2pt]  {$e_1$};
  \draw[rim]   (v2)--(v3) node[midway, above right=-2pt] {$e_3$};
  \draw[rim]   (v1)--(v4) node[midway, below left=-2pt]  {$e_2$};
  \draw[rim]   (v2)--(v4) node[midway, below right=-2pt] {$e_4$};
  \draw[spine] (v1)--(v2) node[midway, above] {$e_0$};

  \node[ver] at (v1) {$1$};
  \node[ver] at (v2) {$2$};
  \node[ver] at (v3) {$3$};
  \node[ver] at (v4) {$4$};

  % the ground
  \draw[line width=0.5pt] (0,-1.40)--(0,-1.58);
  \draw[line width=0.7pt] (-0.24,-1.58)--(0.24,-1.58);
  \draw[line width=0.5pt] (-0.14,-1.69)--(0.14,-1.69);
  \draw[line width=0.5pt] (-0.06,-1.78)--(0.06,-1.78);
  \node[anchor=west, font=\footnotesize] at (0.31,-1.64) {$y_4=0$};

  \node[font=\footnotesize] at (0,-2.24)
       {$c=(2,3,3,3,3)$,\quad $t_e=\tfrac15$};
\end{scope}

% ================= the five incidence vectors =============================
\begin{scope}[shift={(5.70,0.10)}]
  \coordinate (O)  at ( 0,0);
  \coordinate (b1) at (-1.9226, 0.4209);   % (1,-1, 0)
  \coordinate (b2) at (-0.5130,-0.6627);   % (1, 0,-1)
  \coordinate (b3) at (-0.5130, 0.6617);   % (1, 0, 0)
  \coordinate (b4) at ( 1.4096,-1.0836);   % (0, 1,-1)
  \coordinate (b5) at ( 1.4096, 0.2408);   % (0, 1, 0)

  \draw[vec] (O)--(b1) node[left=1pt]         {$b_0$};
  \draw[vec] (O)--(b2) node[below left=-2pt]  {$b_1$};
  \draw[vec] (O)--(b3) node[above left=-2pt]  {$b_2$};
  \draw[vec] (O)--(b4) node[below right=-2pt] {$b_3$};
  \draw[vec] (O)--(b5) node[right=1pt]        {$b_4$};
  \fill (O) circle (1.1pt);

  \node[font=\footnotesize, align=center] at (0,-2.34)
    {$b_0=(1,-1,0)$,\quad $b_1=(1,0,-1)$,\quad $b_2=(1,0,0)$,\\[1.5pt]
     $b_3=(0,1,-1)$,\quad $b_4=(0,1,0)$};
\end{scope}

\end{tikzpicture}

  \caption{The diamond $\dsn_{\Diamond}$.  Left: the graph $K_4$ with the edge
  $34$ deleted and the vertex $4$ grounded, its edges carrying the
  conductances $c=(2,3,3,3,3)$ and the uniform weights $t_e=\tfrac15$;
  line weight is proportional to conductance.  Right: the five reduced
  incidence vectors of \eqref{eq:drops}, pairwise non-proportional, whence
  the five atoms \eqref{eq:diamondatoms} are pairwise non-parallel.}
  \label{fig:diamond}
\end{figure}

\subsection{Volume sampling}\label{sec:averaging}

Volume sampling is the probability measure on $k$-subsets adapted to
\eqref{eq:parseval}, and under it the subset sum dominates on average,
clearing the threshold by a factor of at least $k$ in trace.  Turning
that average into a single subset is the derandomization the problem asks
for.

We work
in the chart of \S\ref{sec:chart}: $v_c=\sqrt{t_c}\,g_c$, the matrix $V$
has rows $v_c\adj$ and satisfies $V\adj V=I_k$, and $P=VV\adj$ by
\eqref{eq:chart}.  For a $k$-subset $C$ write $V_C$ for the $k\times k$
submatrix of $V$ on the rows $C$, so that $P[C]=V_C^{\vphantom{\mathsf{H}}}V_C\adj$.

\begin{lemma}\label{lem:cauchybinet}
Let $X$ be $k\times m$ and $Y$ be $m\times k$.  Then
\[
  \det(XY)\;=\;\sum_{\lvert C\rvert=k}
  \det\bigl(X^{C}\bigr)\det\bigl(Y_C\bigr),
\]
where $X^{C}$ is the submatrix of $X$ on the columns $C$ and $Y_C$ the
submatrix of $Y$ on the rows $C$.
\end{lemma}

\begin{proof}
Write $x_1,\dots,x_m$ for the columns of $X$.  The $j$-th column of $XY$
is $\sum_{c}Y_{cj}x_c$, so multilinearity of the determinant in its
columns gives
\[
  \det(XY)\;=\;\sum_{f}\Bigl(\prod_{j=1}^{k}Y_{f(j)j}\Bigr)
  \det\bigl(x_{f(1)},\dots,x_{f(k)}\bigr),
\]
the sum being over all maps $f:\{1,\dots,k\}\to\{1,\dots,m\}$.  A
non-injective $f$ repeats a column and contributes zero.  An injective $f$
has an image $C$ of size $k$; writing $f=\iota_C\circ\sigma$ with
$\iota_C$ the increasing enumeration of $C$ and $\sigma$ a permutation,
the determinant equals $\operatorname{sgn}(\sigma)\det(X^{C})$, and
summing $\operatorname{sgn}(\sigma)\prod_jY_{\iota_C(\sigma(j))j}$ over
$\sigma$ is the Leibniz formula for $\det(Y_C)$.
\end{proof}

Apply the lemma with $X=V\adj$ and $Y=V$.  The columns of $V\adj$ on $C$
are the conjugated rows of $V$ on $C$, so $(V\adj)^{C}=(V_C)\adj$ and each
term is
$\det\bigl((V_C)\adj\bigr)\det(V_C)=\lvert\det V_C\rvert^{2}
=\det\bigl(V_C^{\vphantom{\mathsf{H}}}V_C\adj\bigr)=\det P[C]$, while the left side
is $\det(V\adj V)=1$.  Setting
\begin{equation}\label{eq:dpp}
  \pi_C\;:=\;\det P[C]
  \;=\;\lvert\det V_C\rvert^{2}
  \;=\;\Bigl(\prod_{c\in C}t_c\Bigr)
       \det\Gram\bigl((g_c)_{c\in C}\bigr)
\end{equation}
--- the last equality because $V_C=\diag(\sqrt{t_c})_{c\in C}M_C$ with
$M_C$ the matrix of the selected vectors --- we conclude that $\pi$ is a
probability measure on $k$-subsets, and that $\pi_C>0$ exactly when
$(g_c)_{c\in C}$ is a basis of $\FF^k$.  Call such a $C$ a \emph{basis} of
the design.  The measure $\pi$ is volume sampling \cite{DRVW2006}; it is
the determinantal measure carried by the bases of the underlying matroid
\cite{Lyons2003}.

\begin{lemma}\label{lem:marginal}
For every atom $c$,
\[
  \Pr\nolimits_{\pi}\bigl[c\in C\bigr]\;=\;\sum_{C\ni c}\pi_C
  \;=\;P_{cc}\;=\;s_c .
\]
\end{lemma}

\begin{proof}
Fix $c$ and let $E_\sigma:=\diag(1,\dots,\sigma,\dots,1)$ carry $\sigma$
in the position $c$.  On one hand
$V\adj E_\sigma V=V\adj V+(\sigma-1)v_c^{\vphantom{\mathsf{H}}}v_c\adj$ is a
rank-one perturbation of the identity, and
$\det(I_k+\lambda uu\adj)=1+\lambda\lVert u\rVert^{2}$, so
\[
  \det\bigl(V\adj E_\sigma V\bigr)\;=\;1+(\sigma-1)\lVert v_c\rVert^{2}
  \;=\;1+(\sigma-1)s_c .
\]
On the other hand Lemma~\ref{lem:cauchybinet} applied to
$X=V\adj E_\sigma$ and $Y=V$ gives, since scaling the column $c$ of
$V\adj$ multiplies every minor containing it by that factor,
\[
  \det\bigl(V\adj E_\sigma V\bigr)
  \;=\;\sum_{\lvert C\rvert=k}\sigma^{[c\in C]}\,\pi_C
  \;=\;\sum_{C\not\ni c}\pi_C\;+\;\sigma\sum_{C\ni c}\pi_C .
\]
Both expressions are affine in $\sigma$; comparing the coefficients of
$\sigma$ gives the claim.
\end{proof}

Summing Lemma~\ref{lem:marginal} over $c$ returns $\sum_cs_c=k$
(Lemma~\ref{lem:trace}), the expected cardinality of $C$, as it must.

\begin{theorem}[Lean]\label{thm:leverageaverage}
For every design over either field,
\[
  \sum_{c}s_c\,g_c^{\vphantom{\mathsf{H}}}g_c\adj\;\psd\;I_k .
\]
\end{theorem}

\begin{proof}
Let $x\ne0$; at $x=0$ both sides vanish.  Cauchy--Schwarz at each atom gives
$\lvert\langle g_c,x\rangle\rvert^{2}\le\ell_c\lVert x\rVert^{2}$, whence
\[
  \lVert x\rVert^{2}\sum_c s_c\lvert\langle g_c,x\rangle\rvert^{2}
  \;=\;\sum_c t_c\,\ell_c\lVert x\rVert^{2}\,
        \lvert\langle g_c,x\rangle\rvert^{2}
  \;\ge\;\sum_c t_c\lvert\langle g_c,x\rangle\rvert^{4} .
\]
Cauchy--Schwarz for the probability weights $t$ bounds the right side
below by
\[
  \Bigl(\sum_ct_c\lvert\langle g_c,x\rangle\rvert^{2}\Bigr)^{2}
  \;=\;\lVert x\rVert^{4}
\]
by \eqref{eq:parseval}.  Dividing by
$\lVert x\rVert^{2}>0$ gives
$x\adj\bigl(\sum_cs_cg_c^{\vphantom{\mathsf{H}}}g_c\adj\bigr)x\ge\lVert x\rVert^2$.
\end{proof}

The leverage-weighted sum is the $\pi$-average of the subset sum.

\begin{theorem}\label{thm:average}
For every design over either field,
\[
  \E_{\pi}\bigl[\SC\bigr]
  \;=\;\sum_{c}s_c\,g_c^{\vphantom{\mathsf{H}}}g_c\adj
  \;\psd\;I_k .
\]
\end{theorem}

\begin{proof}
Exchanging the two summations and applying Lemma~\ref{lem:marginal},
\[
  \E_\pi[\SC]\;=\;\sum_{\lvert C\rvert=k}\pi_C\sum_{c\in C}
  g_c^{\vphantom{\mathsf{H}}}g_c\adj
  \;=\;\sum_{c}\Bigl(\sum_{C\ni c}\pi_C\Bigr)g_c^{\vphantom{\mathsf{H}}}g_c\adj
  \;=\;\sum_c s_c\,g_c^{\vphantom{\mathsf{H}}}g_c\adj ,
\]
and the inequality is Theorem~\ref{thm:leverageaverage}.
\end{proof}

The selection problem is thus the derandomization of
Theorem~\ref{thm:average}: the $\pi$-average of $\SC$ dominates, and one
asks for a single $C$.

\begin{proposition}[Lean]\label{prop:jensen}
$\sum_c t_c\ell_c^{2}\ge k^{2}$, with equality if and only if $\ell_c=k$
for every $c$.
\end{proposition}

\begin{proof}
Cauchy--Schwarz for the weights $t$ gives
$\bigl(\sum_ct_c\ell_c\bigr)^{2}\le\sum_ct_c\ell_c^{2}$, and the left side
is $k^{2}$ by Lemma~\ref{lem:trace}.  Equality holds exactly when $\ell$
is constant, and its $t$-mean is $k$.
\end{proof}

Since $\tr\E_\pi[\SC]=\sum_cs_c\ell_c=\sum_ct_c\ell_c^{2}\ge k^{2}$ while
$\tr I_k=k$, the average clears the threshold by a factor at least $k$ in
trace.  No single subset need exceed it: that is
Corollary~\ref{cor:nomargin}.

Averaging characteristic polynomials rather than matrices gives the
mixture
\begin{equation}\label{eq:mixture}
  q(x)\;:=\;\E_\pi\bigl[\chi_{\SC}(x)\bigr]
  \;=\;\sum_{\lvert C\rvert=k}\pi_C\,\det\bigl(xI_k-\SC\bigr).
\end{equation}
Call a basis $C$ \emph{tight} if $1\in\operatorname{spec}\SC$, and a
design \emph{totally tight} if every basis of it is tight.

\begin{proposition}\label{prop:totaltight}
Every design of Theorem~\ref{thm:tieseverywhere} is totally tight.
\end{proposition}

\begin{proof}
Let $C$ be a $k$-subset.  If $\varphi$ is not injective on $C$ then $\SC$
is singular by Lemma~\ref{lem:parallel}, so $\pi_C=0$ by \eqref{eq:dpp}
and $C$ is not a basis.  Otherwise $\SC=S_{\varphi(C)}$ dominates and does
not dominate strictly, so $\lmin(\SC)=1$ by Lemma~\ref{lem:tiemax}.
\end{proof}

\begin{proposition}[Lean]\label{prop:qonezero}
$q(1)=0$ at every totally tight design.
\end{proposition}

\begin{proof}
$q(1)=\sum_C\pi_C\det(I_k-\SC)$, and every summand with $\pi_C>0$ has $C$
a basis, hence tight, hence $\det(I_k-\SC)=0$.
\end{proof}

\begin{proposition}\label{prop:qone}
The mixture \eqref{eq:mixture} is monic of degree $k$; its subleading
coefficient is
\[
  [x^{k-1}]\,q\;=\;-\,\E_\pi\bigl[\tr\SC\bigr]
  \;=\;-\sum_c t_c\ell_c^{2}\;\le\;-k^{2};
\]
and $q(1)=0$ at every totally tight design.
\end{proposition}

\begin{proof}
Each $\chi_{\SC}$ is monic of degree $k$ and $\sum_C\pi_C=1$, so $q$ is
monic of degree $k$; and $[x^{k-1}]\chi_{\SC}=-\tr\SC$, which gives the
first equality.  The second is Lemma~\ref{lem:marginal} again:
\[
  \E_\pi[\tr\SC]=\sum_C\pi_C\sum_{c\in C}\ell_c
  =\sum_c\ell_c\sum_{C\ni c}\pi_C=\sum_cs_c\ell_c=\sum_ct_c\ell_c^{2},
\]
and the bound is Proposition~\ref{prop:jensen}.  The last assertion is
Proposition~\ref{prop:qonezero}.
\end{proof}

We compute the mixture at the tetrahedron of Example~\ref{ex:tetra},
which by Theorem~\ref{thm:corankone}(iii)--(iv) is the corank-one extremal design at
$k=3$ and uniform weights, and is therefore totally tight by
Proposition~\ref{prop:totaltight}.  Each of its four triples has Gram
matrix $4I_3-J$, with $J$ the all-ones matrix, of determinant $16$; so by
\eqref{eq:dpp}
\[
  \pi_C\;=\;\bigl(\tfrac14\bigr)^{3}\cdot16\;=\;\tfrac14
\]
for each of the four, confirming $\sum_C\pi_C=1$ and giving
$\Pr_\pi[c\in C]=3\cdot\tfrac14=\tfrac34=s_c$ as
Lemma~\ref{lem:marginal} requires.  Every triple has spectrum $\{1,4,4\}$,
so every $\chi_{\SC}(x)=(x-1)(x-4)^{2}$ and the mixture is that same
polynomial:
\begin{equation}\label{eq:qtetra}
  q(x)\;=\;(x-1)(x-4)^{2}\;=\;x^{3}-9x^{2}+24x-16 .
\end{equation}
Its subleading coefficient is
$-9=-\sum_ct_c\ell_c^{2}=-4\cdot\tfrac14\cdot9$, so the bound of
Proposition~\ref{prop:jensen} is attained here; $q(1)=0$, as
Proposition~\ref{prop:qone} requires; and the least root of $q$ sits
at the threshold.  Meanwhile
$\E_\pi[\SC]=\tfrac34\sum_cg_c^{\vphantom{\mathsf{T}}}g_c\tp=3I_3$ while
$\max_C\lmin(\SC)=1$: the average clears the threshold by a factor of
three, and no triple exceeds it.

\begin{remark}\label{rem:interlace}
Only the identities above have been used; three nearby statements do
not follow from them.  First, $q$ is not asserted to
be real-rooted; for a projection determinantal measure that is available
through the strong Rayleigh property
\cite{BorceaBrandenLiggett2009,AnariOveisGharan2015}, and nothing here
needs it.  Second, real-rootedness of $q$ would not imply that some $C$
has $\lmin(\SC)$ at least the least root of $q$: the least root of an
average of real-rooted polynomials may exceed the least root of every
member, and an exact two-polynomial witness is recorded in
Appendix~\ref{app:graveyard}.  Third, the hypothesis under which the
conclusion does follow is a common interlacing of the family
$\{\chi_{\SC}:\pi_C>0\}$, equivalently real-rootedness of every convex
combination of it \cite{MSS2015}; that is a property of the family, to be
verified and not inferred, and it fails for the witness just cited.
Establishing it here, and with it the derandomization of
Theorem~\ref{thm:average}, is Problem~\ref{prob:interlace}; by
Proposition~\ref{prop:qone} any such argument is exact at the extremal
designs, where the least root of $q$ is $1$.
\end{remark}
